\hypertarget{appendix-410a611a5b8d}{}
\section*{C++ modules and reachability}
\section{C++ modules and reachability}\label{c-modules-and-reachability}

\textbf{Curriculum depth:} Introduced · \textbf{First seen:} C007 · \textbf{Latest development:} C007

\subsection{Definition and mechanism}\label{definition-and-mechanism}

A named module imports semantic interfaces rather than textually pasting headers. A BMI supplies compiler-facing interface information; object code still participates in linking. Export visibility is not identical to semantic reachability, and a named-module import is distinct from a macro-bearing header-unit import.

\subsection{Prerequisites}\label{prerequisites}

\begin{itemize}
\tightlist
\item
  \hyperref[appendix-806686359e1b]{Translation units and linking} --- Interface information serves compilation, while object code serves linking; modules change their distribution.
\end{itemize}

\subsection{Used by}\label{used-by}

Later chunks extend this same node; dependent concepts are listed here as they are integrated.

\subsection{Recall and next depth}\label{recall-and-next-depth}

Give an example, explain why the rule exists, and identify which guarantee would be missing if the prerequisite were ignored.

\textbf{Pending:} Deeper ownership, pointer/reference, generic programming and module rules from later sources. This is written coverage, not demonstrated personal mastery.

\subsection{Source and navigation}\label{source-and-navigation}

\href{Continuous\%20Notes/Computing\%20Foundations\%20-\%20From\%20Bits\%20to\%20C++\%20and\%20Systems\%20APIs\#Giving\%20a\%20program\%20a\%20buildable\%20shape}{C007 teaching chapter} · \hyperref[appendix-fd090d982a15]{Source coverage} · \href{}{Dependency Map (02-Foundations-II, appendix-4ad47bdd269a)}
